%Abstract Page

\hbox{\ } \vspace{.7in}
\renewcommand{\baselinestretch}{1}
\small \normalsize

\begin{center}
    \large{{ABSTRACT}}

    \vspace{3em}

\end{center}
\hspace{-.15in}
\begin{tabular}{ll}
    Title of Dissertation:    & {\large  Designing Effective Property-Based Testing Frameworks} \\
    \multicolumn{2}{l}{}                                                                        \\
                              & {\large  Alperen Keleş}                                         \\
                              & {\large Doctor of Philosophy, 2026}                             \\
    \multicolumn{2}{l}{}                                                                        \\
    Dissertation Directed by: & {\large Assistant Professor Leonidas Lampropoulos}              \\
                              & {\large  Department of Computer Science }                       \\
                              & {\large  University of Maryland, College Park }                 \\
\end{tabular}

\vspace{3em}

\renewcommand{\baselinestretch}{2}
\large \normalsize

Property-based testing (PBT) is a widely adopted random testing methodology that
combines user-defined logical propositions with automated test case generation.
%
Over the past two decades, PBT has gained significant traction; today, virtually all
major programming languages host multiple competing PBT frameworks, each innovating in
interface design and underlying techniques.
%
As a result of this proliferation, the PBT landscape of today is significantly
fragmented. Advances in one framework often fail to reach others. Common tips, tricks,
and techniques are rarely evaluated quantitatively. Historical design choices, once
built into popular libraries, can keep the ecosystem from adapting. As a result, the
field still lacks a shared empirical foundation for deciding which PBT designs work
best, when, and why.

This thesis takes a step toward consolidating PBT knowledge. It contributes a
conceptual framework for describing property runners in the wild; a flexible property
language for expressing runners distributed throughout the PBT ecosystem; ETNA, the
first large-scale evaluation and benchmarking platform for PBT; and DIRT, a
domain-specific testing approach that clarifies which challenges remain outside the
reach of general-purpose consolidation. Together, these contributions make it easier to
compare PBT techniques, transfer ideas across frameworks, and build the next generation
of testing tools on shared evidence rather than local convention.
